\begin{figure}[htbp]
	\centering
	\begin{tikzpicture}[font=\small]
		\def\W{4.6}\def\h{0.6}\def\a{0.85}\def\r{0.55}
		% concio
		\draw[trvconcio] (0,-\h) rectangle (\W,\h);
		\draw[trvfibra] (0,-\h-0.12) -- (\W,-\h-0.12);
		% faccia sinistra: azioni con verso opposto a quello positivo
		\draw[trvforza] (0,0) -- ++(-\a,0) node[left] {$N$};
		\draw[trvforza] (0,0) -- ++(0,\a) node[above] {$T$};
		\pic at (0,0) {trvmomgiu};
		\node[below left] at ($(0,0)+(235:0.42)$) {$M$};
		% faccia destra: azioni nei versi positivi
		\draw[trvforza] (\W,0) -- ++(\a,0) node[right] {$N + \Delta N$};
		\draw[trvforza] (\W,0) -- ++(0,-\a) node[right, pos=0.9] {$T + \Delta T$};
		\pic at (\W,0) {trvmomsu};
		\node[above right] at ($(\W,0)+(55:0.42)$) {$M + \Delta M$};
		% risultanti delle azioni distribuite
		\draw[trvforza] ({\W/2},0) -- ++(0.75,0) node[above=1pt, pos=0.6] {$q(z)\,\Delta z$};
		\draw[trvforza] ({\W/2},0) -- ++(0,-0.5) node[right=1pt, pos=0.7] {$p(z)\,\Delta z$};
		\draw[trvcoppia] ({\W/2-0.75},-0.12) arc[start angle=0, end angle=180, radius=0.24];
		\node[below] at ({\W/2-0.99},-0.12) {$c(z)\,\Delta z$};
		\pic at ({\W/2},0) {trvpunto};
		% quota
		\draw[trvrichiamo] (0,-\h-0.2) -- (0,-1.6);
		\draw[trvrichiamo] (\W,-\h-0.2) -- (\W,-1.6);
		\draw[trvquota] (0,-1.45) -- node[below=1pt] {$\Delta z$} (\W,-1.45);
		\node[below left=0pt] at (0,-1.45) {$z$};
		\node[below right=0pt] at (\W,-1.45) {$z + \Delta z$};
	\end{tikzpicture}
	\caption{Concio elementare compreso fra le sezioni di
		ascissa $z$ e $z + \Delta z$ in un tratto regolare.
		Sulla faccia di ascissa $z$ le azioni interne agiscono
		con verso opposto a quello positivo, sulla faccia di
		ascissa $z + \Delta z$ con verso positivo e incrementate
		di $\Delta N$, $\Delta T$, $\Delta M$. Al primo ordine
		in $\Delta z$ le azioni distribuite si riducono alle
		risultanti $p(z)\,\Delta z$, $q(z)\,\Delta z$, $c(z)\,\Delta z$.}
	\label{fig:concio_eq_indefinite}
\end{figure}
